\section{Algorithm}
\label{sec:algorithm}

\subsection{Formal Definition}

Let $G = (V, E)$ be a census-tract adjacency graph with $|V| = n$,
$k$ the number of districts, and $\ec(\pi)$ the edge-cut of plan $\pi$
(the number of edges in $E$ crossing district boundaries).
Equivalently, $1 - \ec(\pi)/|E|$ is monotone in Polsby-Popper: lower
edge-cut corresponds to higher compactness.

\begin{definition}[\shb{}]
\label{def:short-burst}
Given parameters $\ell$ (burst length), $m$ (burst count), $p \in [0, 1]$
(selection percentile), and a base seed $s \in \{0, \ldots, 2^{64}-1\}$:
\begin{enumerate}
\item Compute the initial plan $\pi_0 = \mathrm{METIS}(G, k, s)$.
\item For $i = 0, \ldots, m-1$:
  \begin{enumerate}
  \item Derive the burst seed $s_i = \mathrm{chain\_seed}(s, i)$.
  \item Initialize a \recom{} chain from $\pi_i$ using $s_i$ as the RNG seed.
  \item Run the chain for $\ell$ steps to obtain $\pi_i^{(\ell)}$.
  \item Record the endpoint: $e_i = (\ec(\pi_i^{(\ell)}),\; i,\; \pi_i^{(\ell)})$.
  \item Set $\pi_{i+1} = \pi_i^{(\ell)}$ (restart from endpoint).
  \end{enumerate}
\item Sort endpoints: $e_{(0)} \leq e_{(1)} \leq \cdots \leq e_{(m-1)}$
      by $(\ec, i)$ lexicographically (ascending EC, then burst index as tiebreak).
\item Return $e_{(\lfloor p \cdot m \rfloor)}.\pi$ (the plan at the $p$-th
      quantile of recorded endpoints by edge-cut).
\end{enumerate}
\end{definition}

The default parameters in \texttt{redist-cli} are $\ell = 20$, $m = 50$,
$p = 0.0$ (always return the best endpoint found).

\subsection{Chain-Seed Derivation}

Each burst uses a deterministic, independent RNG seed derived from the
base seed and burst index.
This ensures:
\begin{itemize}
\item \textbf{Reproducibility}: given the same base seed, the same
      sequence of plans is produced on any run.
\item \textbf{Independence}: burst seeds are cryptographically independent,
      preventing correlation between consecutive burst RNG streams.
\item \textbf{No seed reuse}: each burst index maps to a unique seed,
      preventing the same random choices from appearing in different bursts.
\end{itemize}

The derivation uses SHA-256.
Define the domain-separated byte string:
\[
  D(s, i) \;=\;
    \underbrace{\texttt{"SB\_CHAIN\_"}}_{\text{9 bytes}}
    \;\|\;
    \underbrace{i.\mathrm{to\_le\_bytes}()}_{\text{8 bytes}}
    \;\|\;
    \underbrace{\texttt{"\_"}}_{\text{1 byte}}
    \;\|\;
    \underbrace{s.\mathrm{to\_le\_bytes}()}_{\text{8 bytes}},
\]
then:
\[
  \mathrm{chain\_seed}(s, i) = \mathrm{SHA256}(D(s, i))[0..8],
\]
where $[0..8]$ denotes the first 8 bytes of the SHA-256 digest,
interpreted as a little-endian \texttt{u64}.
The prefix serves as domain separator; the implementation uses the
longer string \texttt{"SHORT\_BURST\_CHAIN\_"} for legibility.

This is the same domain-separation pattern used in \be{}~\citep{h1paper}
and \cs{}~\citep{b16paper}: each algorithm uses a unique ASCII prefix
to ensure no seed collision across search modes.

\subsection{Pseudocode}

\begin{algorithm}[H]
\caption{\shb{} Plan Selection}
\label{alg:short-burst}
\begin{algorithmic}[1]
\Require Graph $G$, districts $k$, burst length $\ell$, burst count $m$,
         percentile $p$, base seed $s$
\Ensure A redistricting plan $\pi^*$
\State $\pi_0 \gets \mathrm{METIS}(G, k, s)$
\State $\mathit{endpoints} \gets []$
\For{$i \gets 0$ \textbf{to} $m-1$}
  \State $s_i \gets \mathrm{chain\_seed}(s, i)$
  \State $\mathit{rng} \gets \mathrm{SmallRng::seed\_from\_u64}(s_i)$
  \State $\mathit{chain} \gets \mathrm{RecomChain}(\pi_i,\; G,\; k)$
  \For{$j \gets 0$ \textbf{to} $\ell - 1$}
    \State $\mathit{chain}.\mathrm{step}(\mathit{rng})$
  \EndFor
  \State $\pi_{i+1} \gets \mathit{chain}.\mathrm{assignment}$
  \Comment{endpoint, not minimum}
  \State $\mathit{endpoints}.\mathrm{push}\!\left((\ec(\pi_{i+1}),\; i,\; \pi_{i+1})\right)$
\EndFor
\State $\mathit{endpoints}.\mathrm{sort\_by}\!\left(\ec \text{ asc},\; i \text{ asc}\right)$
\Comment{stable, deterministic}
\State \Return $\mathit{endpoints}[\lfloor p \cdot m \rfloor].\pi$
\end{algorithmic}
\end{algorithm}

\subsection{CLI Invocation}

\begin{verbatim}
redist state --state NC --year 2020
    --search short-burst
    --burst-length 20
    --n-bursts 50
    --percentile 0.0
\end{verbatim}

The percentile parameter \texttt{--percentile 0.0} selects the best
endpoint (minimum edge-cut) across all bursts.
Setting \texttt{--percentile 0.5} returns the median endpoint, which
may be preferable when statistical properties of the neighbourhood are
of interest rather than the optimum.

\subsection{Stochastic Dominance}

\begin{theorem}[Short-Burst Stochastic Dominance]
\label{thm:dominance}
Let $\pi^{\shb}$ be the plan returned by \shb{} with $\ell$ steps per
burst and $m$ bursts, and let $\pi^{\recom}$ be the plan at the end
of a single \recom{} chain of $\ell \cdot m$ steps (same total step
budget), both starting from the same METIS initialisation.
Then for any threshold $c > 0$:
\[
  \Pr\!\left[\ec(\pi^{\shb}) \leq c\right]
  \;\geq\;
  \Pr\!\left[\ec(\pi^{\recom}) \leq c\right].
\]
That is, \shb{} stochastically dominates single-chain \recom{} of the
same total length for finding low-edge-cut plans.
\end{theorem}

\begin{proof}[Proof sketch]
The key observation is that after burst $i$, the chain restarts from
$\pi_{i+1}$, which has already been selected as the best-available
starting point for the next burst.
By contrast, the single chain at step $i \cdot \ell$ has no such
restart---it may be at a higher-edge-cut state than any endpoint
recorded so far.

Formally, let $\mathcal{E}_i = \min_{j \leq i} \ec(\pi_j^{(\ell)})$
be the best edge-cut seen after $i$ bursts.
For the single chain, let $\mathcal{E}'_T = \ec(\pi'_T)$ be the
edge-cut at step $T = i \cdot \ell$.

After $i$ bursts, the \shb{} chain is at state $\pi_{i+1}$ with
$\ec(\pi_{i+1}) \geq \mathcal{E}_i$ (the restart is from the endpoint,
not the best).
However, the next burst starts from $\pi_{i+1}$, which is constrained
to have $\ec(\pi_{i+1}) \leq \ec(\pi_i^{(\ell)}) + \Delta$ for some
bounded drift $\Delta$ over $\ell$ steps.

By induction on the number of bursts: the restart mechanism ensures
that the chain cannot drift far from the best-found plan.
The $p = 0.0$ selection then returns the minimum over all burst
endpoints, which is at least as good as any single endpoint of the
same-budget chain.
\end{proof}

\begin{remark}
This is a heuristic dominance claim: the proof sketch above bounds the
drift per burst but does not formally characterise the mixing time of
the local neighbourhood.
A formal proof would require mixing time bounds on the restriction of
the \recom{} chain to the $\delta$-neighbourhood of a good plan---a
technically involved argument beyond the scope of this paper.
The empirical evidence in Section~\ref{sec:comparison} provides strong
evidence for the dominance claim in practice.
\end{remark}

\subsection{Why Endpoint Retention is Non-Negotiable}

We emphasise the correctness requirement at the heart of Short-Burst.

\paragraph{Breaking the Markov property.}
If one were to restart from the within-burst minimum rather than the
endpoint, the resulting sequence would not be a Markov chain on $\mathcal{P}$.
The within-burst minimum depends on the entire trajectory of the burst,
not just the current state.
A chain that restarts from trajectory-dependent states violates the
Markov property: the future evolution depends on past history beyond
the current state.

\paragraph{Systematic bias.}
Selecting the within-burst minimum introduces systematic bias: the
selected plan is systematically lower in edge-cut than the chain's
stationary distribution would imply.
This is not merely a theoretical concern---it means the selected plan
cannot be interpreted as a sample from the stationary distribution,
and the sequence of selected plans does not approximate any well-defined
distribution on $\mathcal{P}$.

\paragraph{Correct behaviour.}
Endpoint retention preserves both properties.
The endpoint is a valid state of the \recom{} chain---the chain ended
there, and the next chain starts there, maintaining the Markov property.
The sequence of endpoints is a valid stochastic process on $\mathcal{P}$,
and the selection of the minimum endpoint is a well-defined post-processing
step that does not affect the validity of each individual restart.
